\section{A Laplace lemma with a parameter}\label{app:laplace}

This appendix proves the Laplace lemma that Appendix~\ref{app:torus} uses
for Theorem~\ref{thm:local}. On $\mathbb T^m=(\R/2\pi\Z)^m$ we use
$\lVert\varphi\rVert=\max_i\operatorname{dist}(\varphi_i,2\pi\Z)$. For
$x\in\R^m$ we write $\lVert x\rVert=\max_i|x_i|$ and $|x|$ for the Euclidean
norm. Repeated indices are summed from $1$ to $m$.

\begin{lemma}[Laplace integrals with a parameter]\label{lem:laplace}
Let $\mathcal K$ be a compact metric space and $m\ge1$. Let
$\Psi,a:\mathcal K\times\mathbb T^m\to\mathbb C$ be such that the
$\varphi$-derivatives of $\Psi$ up to order $4$ and of $a$ up to order $2$
exist and are jointly continuous on $\mathcal K\times\mathbb T^m$. Assume that
for every $\alpha\in\mathcal K$
\begin{enumerate}[(i)]
\item $\Psi(\alpha,0)\in\R$ and $\nabla_\varphi\Psi(\alpha,0)=0$,
\item $\mathcal J(\alpha)=-\nabla^2_\varphi\Psi(\alpha,0)$ is real, symmetric
and positive definite, and
\item $\operatorname{Re}\Psi(\alpha,\varphi)<\Psi(\alpha,0)$ for every
$\varphi\ne0$ in $\mathbb T^m$.
\end{enumerate}
Then, uniformly for $\alpha\in\mathcal K$ and $T\ge1$,
\[
\int_{\mathbb T^m}e^{T\Psi(\alpha,\varphi)}a(\alpha,\varphi)\,d\varphi
=e^{T\Psi(\alpha,0)}\Bigl(\frac{2\pi}T\Bigr)^{m/2}\det\mathcal J(\alpha)^{-1/2}
\Bigl(a(\alpha,0)+O\Bigl(\frac1T\Bigr)\Bigr).
\]
\end{lemma}

\begin{proof}
By compactness
$B_\Psi=\max_{|\beta|\le4}\sup|\partial^\beta_\varphi\Psi|$ and
$B_a=\max_{|\beta|\le2}\sup|\partial^\beta_\varphi a|$ are finite, the least
eigenvalue of $\mathcal J(\alpha)$ has a positive minimum $h_{\min}$ over
$\mathcal K$, and $\det\mathcal J(\alpha)\le(mB_\Psi)^m$. Below, $C$ depends
only on $m$, $h_{\min}$, $B_\Psi$, $B_a$ and the numbers $\varepsilon$ and
$\eta$ chosen below, and it may change from line to line. So every bound
below is uniform on $\mathcal K$.

\emph{Localization.} Choose $\varepsilon\in(0,1)$ with
$m^3B_\Psi\varepsilon\le\frac32h_{\min}$. The set of $(\alpha,\varphi)$ with
$\lVert\varphi\rVert\ge\varepsilon$ is compact, and by (iii) the continuous
function $\operatorname{Re}\Psi(\alpha,\varphi)-\Psi(\alpha,0)$ is negative on
it. So it is at most $-\eta$ there for some $\eta>0$, and the integral over
this set has modulus at most
$(2\pi)^mB_ae^{T(\Psi(\alpha,0)-\eta)}\le CT^{-1-m/2}e^{T\Psi(\alpha,0)}$.

\emph{Scaling.} Reduction modulo $2\pi$ identifies the cube
$\lVert x\rVert\le\varepsilon$ in $\R^m$ with the set
$\lVert\varphi\rVert\le\varepsilon$, and it preserves measure. Let
$R(\alpha,\varphi)=\Psi(\alpha,\varphi)-\Psi(\alpha,0)
+\frac12\varphi^\top\mathcal J(\alpha)\varphi$. By (i) and (ii), $R$ and its
derivatives of order $1$ and $2$ vanish at $\varphi=0$. Taylor's theorem with
integral remainder, which holds for complex values, gives on the cube
$R=\frac16\Psi_{ijk}\varphi^i\varphi^j\varphi^k+\omega_0$ with
$|\omega_0|\le C\lVert\varphi\rVert^4$, where $\Psi_{ijk}$ are the third
derivatives at $(\alpha,0)$. It also gives
$|R|\le\frac{m^3}6B_\Psi\lVert\varphi\rVert^3\le\frac14h_{\min}|\varphi|^2$,
since $\sum_{|\beta|=3}1/\beta!=m^3/6$ and
$\lVert\varphi\rVert^3\le\varepsilon|\varphi|^2$. We substitute
$\varphi=x/\sqrt T$ and put $w(x)=TR(\alpha,x/\sqrt T)$. The integral over the
cube is
\[
T^{-m/2}e^{T\Psi(\alpha,0)}\int_{\lVert x\rVert\le\varepsilon\sqrt T}
e^{-\frac12x^\top\mathcal Jx}\,e^{w(x)}a(\alpha,x/\sqrt T)\,dx,
\]
and $|w(x)|\le\frac14h_{\min}|x|^2$.

\emph{Taylor expansion.} Let $P(x)=\frac16\Psi_{ijk}x^ix^jx^k$,
$a_0=a(\alpha,0)$ and $a_i=\partial_{\varphi_i}a(\alpha,0)$. Then
$w=T^{-1/2}P+\omega_1$ with $|\omega_1|\le CT^{-1}|x|^4$, and
$|w|\le CT^{-1/2}|x|^3$. Since $|e^w-1-w|\le\frac12|w|^2e^{|w|}$, this gives
$|e^w-1-T^{-1/2}P|\le CT^{-1}(|x|^4+|x|^6)e^{h_{\min}|x|^2/4}$. Taylor's
theorem for $a$ to orders $1$ and $2$ gives
$|a(\alpha,x/\sqrt T)-a_0|\le CT^{-1/2}|x|$ and
$|a(\alpha,x/\sqrt T)-a_0-T^{-1/2}a_ix^i|\le CT^{-1}|x|^2$. Multiplying out,
\[
e^{w}a(\alpha,x/\sqrt T)=a_0+T^{-1/2}\bigl(a_0P(x)+a_ix^i\bigr)+\omega(x),
\qquad|\omega(x)|\le CT^{-1}(1+|x|)^6e^{h_{\min}|x|^2/4}.
\]

\emph{Integration.} Since $x^\top\mathcal Jx\ge h_{\min}|x|^2$, the term
$\omega$ contributes at most $CT^{-1}$ to the integral, because
$(1+|x|)^6e^{-h_{\min}|x|^2/4}$ is integrable. Outside the cube
$|x|>\varepsilon\sqrt T$, so extending the integrals of the other 2 terms to
$\R^m$ changes them by at most $Ce^{-h_{\min}\varepsilon^2T/4}\le CT^{-1}$. On
$\R^m$ we have
$\int e^{-\frac12x^\top\mathcal Jx}dx=(2\pi)^{m/2}\det\mathcal J^{-1/2}$, and
the odd polynomial $a_0P(x)+a_ix^i$ integrates to $0$ against this weight.
This is why the error
is $O(1/T)$ and not $O(T^{-1/2})$. So the integral over the cube is
$e^{T\Psi(\alpha,0)}T^{-m/2}\bigl[(2\pi)^{m/2}\det\mathcal J^{-1/2}a_0
+O(T^{-1})\bigr]$. Since $\det\mathcal J\le(mB_\Psi)^m$, this error and the
one from the localization are $O(1/T)$ relative to
$e^{T\Psi(\alpha,0)}(2\pi/T)^{m/2}\det\mathcal J^{-1/2}$.
\end{proof}
